\chapter{Motion Sickness}

\section*{Definition}
A syndrome of nausea, dizziness, and autonomic symptoms triggered by real or perceived motion, arising from sensory conflict between visual, vestibular, and proprioceptive inputs.

\section*{When to See a Doctor}
See your doctor if:
\begin{itemize}
  \item Motion sickness with severe or persistent vomiting causing dehydration, or with neurological symptoms
\end{itemize}

\section*{How It Develops}
The sensory conflict theory posits motion sickness occurs when information from the vestibular system (detecting motion) does not match visual and proprioceptive inputs. Autonomic activation via the brainstem vomiting center leads to nausea, pallor, cold sweats, and vomiting. Central habituation reduces symptoms over repeated exposure.

\section*{Causes}
\begin{itemize}
  \item Car travel (especially winding roads)
  \item Boat or sea travel (sea sickness)
  \item Airplane travel
  \item Virtual reality or video simulations
  \item Amusement park rides
  \item Migraine-associated vertigo
  \item Inner ear disorders (Meniere disease, labyrinthitis)
  \item Pregnancy (first trimester)
  \item Migraine headache
  \item Panic disorder or anxiety
\end{itemize}

\section*{Related Symptoms}
\begin{itemize}
  \item Nausea
  \item Dizziness
  \item Sweating
  \item Headache
  \item Pallor
\end{itemize}

\section*{Medical Evaluation}
\begin{itemize}
  \item Diagnosis is primarily clinical based on the classic history of nausea, dizziness, and pallor triggered by motion (car, boat, plane, VR) with relief when the motion stops.
  \item Neurological examination including nystagmus assessment and Dix-Hallpike manoeuvre to differentiate from vestibular disorders (BPPV, labyrinthitis).
  \item No routine laboratory or imaging tests are required; the diagnosis is made on history alone in typical cases.
  \item Referral to an ENT specialist or neurologist is indicated if symptoms occur without motion exposure, suggesting an underlying vestibular or neurological disorder.
\end{itemize}

\section*{Treatment Options}
\begin{itemize}
  \item Preventive pharmacological therapy includes antihistamines (dimenhydrinate, meclizine, cyclizine) taken 30-60 minutes before travel, or scopolamine transdermal patch applied 4-6 hours before departure.
  \item Acute symptoms are managed with promethazine or ondansetron for nausea; lying down with eyes closed and focusing on a fixed horizon point helps reduce sensory conflict.
  \item Vestibular habituation exercises (Cawthorne-Cooksey exercises) can reduce susceptibility over time through repeated graded exposure to motion stimuli.
  \item For severe or refractory motion sickness, consider a combination of scopolamine and an antihistamine under medical guidance.
\end{itemize}

\section*{Self-Care and Home Management}
\begin{itemize}
  \item Position yourself where motion is least felt — in a car, sit in the front passenger seat; on a boat, stay on deck and focus on the horizon; on a plane, sit over the wings.
  \item Avoid reading, using screens, or focusing on nearby objects during travel, as visual-vestibular conflict is the core mechanism of motion sickness.
  \item Take prophylactic medication 30-60 minutes before travel: antihistamines (dimenhydrinate, meclizine) or scopolamine patches for longer journeys.
  \item Use acupressure bands (Sea-Bands) on the P6 (Neiguan) point on the inner wrist, which may provide mild relief through neuromodulation.
\end{itemize}

\section*{References}

\begin{thebibliography}{99}
\bibitem{motion_sickness:harrison-motion} Longo DL, et al. \emph{Harrison's Principles of Internal Medicine}. 21st ed. McGraw-Hill; 2022. Chapter 31: Motion Sickness.
\bibitem{motion_sickness:uptodate-motion} Hain TC, et al. Motion sickness. In: UpToDate, Post TW (Ed), UpToDate, Waltham, MA; 2025.
\bibitem{motion_sickness:koch} Koch A, et al. Motion sickness: a clinical review. \emph{JAMA}. 2023;330(1):59-70.
\bibitem{motion_sickness:golding} Golding JF, et al. Motion sickness: pathophysiology and prevention. \emph{Auton Neurosci}. 2024;248:103138.
\bibitem{motion_sickness:shupak} Shupak A, et al. The vestibular system and motion sickness. \emph{N Engl J Med}. 2023;388(22):2078-2088.
\bibitem{motion_sickness:spinks} Spinks AB, et al. Scopolamine for the prevention of motion sickness. \emph{Cochrane Database Syst Rev}. 2023;3:CD002851.
\end{thebibliography}
